% ============================================================================
% SEC05.TEX - Section 6.5: Raycast
% ============================================================================

\section{Raycast}

\begin{learningoutcomes}
\item Memahami konsep Raycast
\item Menggunakan Physics.Raycast
\item Mengimplementasikan line of sight detection
\item Menggunakan Raycast untuk shooting mechanics
\end{learningoutcomes}

\subsection{Pengertian Raycast}

\begin{definisi}[Raycast]
Raycast adalah teknik untuk mendeteksi objek dengan menembakkan sinar imajiner dari satu titik ke arah tertentu \cite{UnityPhysics2024}.
\end{definisi}

\subsection{Contoh Implementasi}

\begin{lstlisting}[language={[Sharp]C}]
public class RaycastExample : MonoBehaviour
{
    void Update()
    {
        // Simple raycast
        if (Input.GetMouseButtonDown(0))
        {
            Ray ray = Camera.main.ScreenPointToRay(Input.mousePosition);
            RaycastHit hit;
            
            if (Physics.Raycast(ray, out hit))
            {
                Debug.Log("Hit: " + hit.collider.name);
                Debug.Log("Distance: " + hit.distance);
                Debug.Log("Point: " + hit.point);
            }
        }
        
        // Raycast from object
        RaycastHit hitInfo;
        if (Physics.Raycast(transform.position, transform.forward, out hitInfo, 10f))
        {
            if (hitInfo.collider.CompareTag("Enemy"))
            {
                AttackEnemy(hitInfo.collider.gameObject);
            }
        }
    }
}
\end{lstlisting}

\begin{catatan}
Raycast sangat berguna untuk shooting, line of sight, dan interaction detection. Gunakan LayerMask untuk filter objek yang diraycast.
\end{catatan}
